\section{Stochastic Optimal Control: Formulation and Audit}
\label{ch:optimalcontrol}

The optimal-control extension is the point at which statistical approximation, numerical transcription, and physical modeling become inseparable. A chance constraint is evaluated on a residual produced by a trajectory, but the trajectory itself is generated by a discretized differential equation. Consequently, a small test violation can coexist with a failed nonlinear program, and a reported objective can depend on the quadrature convention rather than on the physical control effort. This chapter states the formulation used by the local implementation, identifies the differences from the published lunar benchmark, and derives an audit identity that limits the interpretation of the local costs.

\subsection{Problem Formulation and Residual Constraints}

Let $y(t)\in\mathbb R^{n_y}$ and $u(t)\in\mathbb R^{n_u}$ denote state and control trajectories, and let $\xi$ denote uncertain parameters. A generic finite-horizon problem is
\begin{equation}
  \begin{aligned}
    \minimize_{y,u}\quad &J(y,u)=\int_0^{t_f}L(y(t),u(t))\,dt+\Psi(y(t_f)),\\
    \text{subject to}\quad &\dot y(t)=F(y(t),u(t),\xi),\\
    &y(0)\in\mathcal Y_0,\quad u(t)\in\mathcal U,\\
    &\Prob(g_\ell(y,u,\xi)\leq0)\geq1-\varepsilon_\ell,
    \quad \ell=1,\ldots,m.
  \end{aligned}
  \label{eq:ocp-general}
\end{equation}
The residual $g_\ell$ may be a terminal event, a path event evaluated at a time $t$, or a vector event reduced to a scalar statistic. The distinction matters: a terminal probability is not a pathwise probability over a continuum of times. In the implementation, path constraints are evaluated at mesh nodes. They are therefore nodewise surrogate constraints unless an additional union or supremum construction is imposed.

For training samples $\{\xi_i\}_{i=1}^{N}$ and a decision vector $z$ containing the nodal states and controls, the residual observations are
\begin{equation}
 r_{\ell i}(z)=g_\ell(y(z),u(z),\xi_i).
\end{equation}
As in Chapters~4--6, the empirical indicator constraint is replaced by an integrated-kernel average, here written in violation form,
\begin{equation}
\hat v_{\ell,h}(z)=\frac{1}{N}\sum_{i=1}^{N}
 \overline K\!\left(\frac{r_{\ell i}(z)}{h}\right),
 \qquad \hat v_{\ell,h}(z)\leq\varepsilon_\ell,
 \label{eq:ocp-kde}
\end{equation}
where $K$ has CDF $F_K$ and $\overline K(t)=1-F_K(-t)$ is the integrated
violation kernel. Hence $\hat v_{\ell,h}=1-\hat p_h$ for the residual
$g_\ell$, in the notation of Chapter~4. Thus positive residuals contribute near
one and negative residuals contribute near zero. Eq.~\eqref{eq:ocp-kde} is a
deterministic NLP constraint after the samples and bandwidth have been fixed.
It does not by itself provide a confidence level for the unknown probability.
The biased construction discussed in Chapter~5 gives a fixed-decision
pointwise envelope; the local biased KDE is recorded as an
implementation surrogate and is not identified with the Split--Bernstein
construction of \citet{keil2020,keil2021}.

\subsection{Discretization and Continuation}

The local fallback solver uses $n_t$ uniformly spaced nodes $0=t_0<\cdots<t_{n_t-1}=t_f$ and forward Euler defects. With $\Delta t_j=t_{j+1}-t_j$,
\begin{equation}
 y_{j+1}-y_j-\Delta t_j F(y_j,u_j,\xi_0)=0,
 \qquad j=0,\ldots,n_t-2,
 \label{eq:euler-defect}
\end{equation}
where $\xi_0$ is absent from the deterministic dynamics in the lunar proxy. The collocation vector includes all nodal controls, although the last control does not enter the forward Euler defects. Equality constraints enforce the initial state, the defects, and terminal velocity. Terminal position remains a chance-constrained quantity. Inequality constraints enforce control bounds and the KDE approximations of the terminal and path residuals.

The implementation separates four operations. The problem module defines dynamics and residuals; the collocation module maps a vector $z$ to nodal states and controls; the KDE module evaluates integrated-kernel statistics; and the solver module supplies the sequential least-squares quadratic programming (SLSQP) solver with equality and inequality callables. This separation makes it possible to inspect a statistical estimator without silently changing the dynamics. It also exposes a numerical limitation: the same nodal control can contribute to a path residual and to the objective while not contributing to the state defects when it is the final control. That final control is therefore an explicit transcription defect, not evidence about physical fuel use. The limitation is quantified below rather than hidden in a single objective number.

The continuation sequence is nominal, worst-case, scenario, unbiased KDE, and biased KDE. A stage is warm-started with the last successful vector. Warm starts are consistent with the continuation strategy described by \citet{keil2021warm}, but the local implementation does not use Legendre--Gauss--Radau (LGR) nodes, adaptive mesh refinement, IPOPT, or GPOPS-II. A failed SLSQP stage is retained in the result ledger with its status and constraint margin; it is not promoted to a feasible solution because its held-out violation happens to be below the target.

\subsection{Local Lunar Benchmark}
\label{sec:lunar-benchmark}

The local benchmark has altitude $y_1$, vertical velocity $y_2$, and nonnegative thrust $u$:
\begin{equation}
 \dot y_1=y_2,\qquad \dot y_2=-g_{\mathrm M}+u,\qquad g_{\mathrm M}=1.622,\qquad 0\leq u\leq3.
 \label{eq:lunar-dynamics}
\end{equation}
The local horizon is fixed at $t_f=20$, the initial state is $(y_1(0),y_2(0))=(100,0)$, and only terminal velocity is an equality. Terminal altitude is subject to
\begin{equation}
 \Prob\left(\left|y_1(t_f)-\xi_1\right|-\eta>0\right)\leq\varepsilon_a,
 \quad \xi_1\sim\mathcal N(0,0.1^2),\quad \eta=0.25,\quad \varepsilon_a=0.1.
 \label{eq:lunar-terminal-risk}
\end{equation}
The nodewise path event is
\begin{equation}
 \Prob\left(u(t_j)+\xi_2-3>0\right)\leq\varepsilon_b,
 \qquad \varepsilon_b=0.01,
 \label{eq:lunar-path-risk}
\end{equation}
with the locally documented mixture $\xi_2\sim0.6\mathcal N(-0.25,0.10^2)+0.4\mathcal N(0.35,0.15^2)$. The mixture is a synthetic stress test for an asymmetric bimodal disturbance; it is not a claim that the local data reproduce the distribution used in the published study.

The comparison with Keil et al. must be made at the level of formulation. Their lunar study uses $y_1(0)=10$, $y_2(0)=-2$, free final time, free terminal altitude, terminal velocity zero, and $J=\int u\,dt$; it solves the resulting problem with LGR/GPOPS-II/SNOPT and reports costs around $9.09$ for the chance-constrained methods \citep{keil2020,keil2021}. The local problem fixes a different initial altitude, initial velocity, horizon, mesh, and solver. It also uses an 11-node forward-Euler transcription solved by SLSQP, declared synthetic normal and two-mode mixture laws in place of the published uncertainty model, and a trapezoidal nodal objective rather than the physical integral $\int u\,dt$ (\cref{sec:quadrature-audit}). The two objective values therefore cannot be compared numerically. The local experiment is a pipeline and audit experiment, not a replication.

\subsection{Conservation and Quadrature Audit}
\label{sec:quadrature-audit}

The local transcription provides a stronger diagnostic than a solver status. Summing the velocity defects over the $n_t-1$ intervals gives
\begin{equation}
  y_{2,n_t-1}-y_{2,0}=\sum_{j=0}^{n_t-2}\Delta t_j(-g_{\mathrm M}+u_j).
 \label{eq:velocity-telescope}
\end{equation}
With uniform $\Delta t=2$, $y_{2,0}=y_{2,n_t-1}=0$, and $g_{\mathrm M}=1.622$, every feasible trajectory satisfies
\begin{equation}
 \sum_{j=0}^{n_t-2}\Delta t\, u_j=g_{\mathrm M} t_f=32.44.
 \label{eq:fuel-invariant}
\end{equation}
Thus, in this fixed-horizon model, the physically consistent piecewise-constant fuel integral is invariant over the feasible set. Trajectory shape still matters because terminal altitude imposes the weighted moment
\begin{equation}
 \int_0^{t_f}(t_f-t)u(t)\,dt
 =\frac{g_{\mathrm M} t_f^2}{2}-y_1(0)=224.4,
 \label{eq:altitude-moment}
\end{equation}
when terminal altitude is zero. The stochastic terminal event changes the admissible timing of thrust, not the invariant integral implied by terminal velocity.

The code records a trapezoidal integral over all $n_t$ controls, whereas the Euler defects use only $u_0,\ldots,u_{n_t-2}$. On a uniform mesh,
\begin{equation}
 J_{\mathrm{trap}}=\Delta t\left(\frac{u_0}{2}+\sum_{j=1}^{n_t-2}u_j+\frac{u_{n_t-1}}{2}\right)
 =32.44+\frac{\Delta t}{2}(u_{n_t-1}-u_0).
 \label{eq:quadrature-mismatch}
\end{equation}
The nominal record $J=29.8052$ is explained by the endpoint term: the stored solution has approximately $u_0=2.6348$ and $u_{n_t-1}=0$. The later records at $32.44$ are not evidence that robust control requires more physical fuel. They reveal that the endpoint contribution becomes different when the optimizer is warm-started under another constraint. This audit is a substantive result of the local implementation and is the reason Chapter~9 reports the lunar objective only as a transcription diagnostic.

\subsection{Interpretation and Transfer Limits}

This section interprets the lunar stage ledger and states what the local experiment can and cannot transfer from the anchor papers. It first applies the solver-feasibility and validation gates to the continuation stages, then states the transfer hypotheses that connect the local run to the anchor papers, and finally delimits the claims that the formulation supports.

\subsubsection{Stage feasibility gates}

The stage ledger is interpreted by two gates: solver feasibility and independent validation. The nominal stage succeeds but has terminal test violation $0.51425$. The worst-case, scenario, unbiased KDE, and biased KDE stages report SLSQP failure; their controls remain diagnostic records, not feasible solutions. Their negative chance margins (approximately $-0.738$, $-0.280$, $-0.0209$, and $-0.0867$) cannot be repaired by a low terminal frequency, even where it lies below $0.1$ (scenario and biased KDE stages), especially because no pathwise held-out validation was performed. The local result is therefore evidence about numerical difficulty and residual statistics, not a successful robust landing.

Warm starts can reduce discontinuities between continuation stages, but they cannot create feasibility when the fixed horizon, transcription, and risk margins are incompatible. The continuation of \citet{keil2021warm} couples kernel changes to increasing samples and mesh refinement; the local sequence changes the estimator at a fixed coarse mesh. It is thus a controlled implementation analogue, not a replication of the published convergence workflow.

\subsubsection{Keil and Schuster transfer hypotheses}

The optimal-control experiment is most informative when stated as three hypotheses rather than as a search for a single best number. First, the Keil hypothesis is that the one-sided Epanechnikov construction with $b(h)=h$ should not underestimate the empirical violation of the residual samples used in the NLP. This is a pointwise property of the integrated kernel. It does not imply that the returned trajectory is feasible when the NLP fails, and it does not provide a finite-sample statement about the unknown distribution. Second, the Schuster hypothesis is that a Gaussian residual KDE with a shrinking bandwidth and increasing sample size should approach the true probability functional when the required convergence assumptions hold. This is a sequence statement, not a claim about the $N=250$ run. Third, the physical hypothesis is that a consistent discretization should make the reported control objective agree with the integral implied by the dynamics. The fixed-horizon identity above provides a direct falsification test for this hypothesis.

The local model satisfies only fragments of these transfer conditions. The biased KDE residual is scalar and uses the same $b(h)=h$ translation as the Keil construction, so the kernel inequality can be checked numerically. The samples are nevertheless IID draws from declared synthetic laws, not the MCMC chains used in the optimal-control paper. The local solver has no mesh-error estimate and no adaptive LGR refinement. The Schuster setting is also only partially represented: the residual KDE is scalar, the sample size is finite, and the static chance boundary is active. Uniform convergence over the decision set is not measured. In addition, the static objective is not globally monotone on its box and the chance-feasible geometry is nonlinear, so the objective and regularity assumptions in the convergence theory cannot be silently assumed.

Adding the Keil warm-start idea to the three hypotheses above gives four transfer claims. The biased-kernel upper-bound claim is represented locally by $b(h)=h$ and pointwise scalar checks; a lower biased KDE violation is evidence of the deliberate one-sided shift, not reproduction of Keil's MCMC/LGR method, and the claim still requires the stated sampling conditions and a feasible nonlinear-program solution. The Schuster convergence claim is represented only in the static and dispatch paper-aligned profiles of Chapter~9, by nested $N$ and shrinking $h_N$ diagnostics; a stable paper-aligned sequence there is numerical evidence for those models, not a uniform-convergence proof, and the claim still requires uniform convergence and a limiting optimizer. The Keil warm-start idea is represented in the lunar run by warm-started stages that change the estimator at a fixed bandwidth, and bandwidth continuation appears only in the same paper-aligned profiles; the published MCMC schedule, mesh error control, and LGR refinement remain absent. The physical-fuel comparison is represented by the telescoping defect audit, in which a changing lunar objective with invariant interval fuel identifies a transcription defect; consistent quadrature and model equivalence are still required before the objectives can be compared. The experiment therefore diagnoses the estimator--trajectory chain without collapsing statistical and physical errors into one metric.

\subsubsection{Scope of valid claims}

The formulation is adequate for a controlled comparison of residual estimators in low dimension. It is not adequate for a claim about continuous-time path safety, because nodewise constraints do not bound the inter-node maximum. It is not adequate for a physical fuel comparison until the quadrature is made consistent with the defect control parameterization. Finally, the finite-dimensional KDE constraint inherits the statistical qualifications developed earlier: convergence of a KDE sequence under assumptions, as in \citet{schuster2025}, does not give a finite-sample confidence statement for the particular 250-sample local run. These limits determine how the lunar records are interpreted in Chapters~9 and~10.
